\section{Implementation}
\subsection{Runtime integration}
DataVine is implemented within the TaskVine codebase using a C workflow
runtime, scheduler, controller protocol, and worker agent, with a Python
workflow interface. One workflow process owns a reactor and the sole manager
API owner. Static graph submission declares task and edge bounds; dynamic
extension retains the same data-identity and attempt contracts. Python
functions use a fork-based execution library. Each logical node generates an
individual physical attempt in the failure-free experiments; retries remain
visible rather than being hidden as task fusion.

The baseline service demand is therefore mixed: one manager reactor drains
dependency and dispatch events while also receiving the data-lifecycle events
caused by intermediate objects. Replica or checkpoint decisions, stage-in and
stage-out acknowledgements, and object-retention or garbage-collection updates
all consume that serialized control path, even when the payload itself is
transferred peer-to-peer. The cost grows with the number of data items and
their consumers, so a worker-side compute shortage is not required for the
manager to saturate.

The service keeps the workflow manager's reactor as the serialized owner of
logical scheduling, while the data-plane endpoint has independent RPC event
threads and Controller persistence workers. RPC threads process authenticated
agent metadata and object requests; persistence workers perform the queued
stream, fsync, verification, and commit path. The default service uses one RPC
thread and 16 persistence workers. The separation is configurable for the
scaling validation, but shared metadata state still imposes a measured limit.

The implementation uses explicit generated-input descriptors and data IDs.
The worker agent handles local materialization and source resolution; the
logical scheduler need not insert a concrete peer endpoint into every task
at placement time. Controller requests are authenticated, and the wire
contracts are versioned. Small callable payloads may be inlined within the
existing size bound, while larger payloads use the corresponding source path.
These encoding optimizations reduce some setup work but are not themselves
the contribution evaluated here.

\subsection{Preparation ablation and observability}
We added a worker environment control selecting deferred or eager preparation.
Deferred is the unchanged default. Eager skips only the check that postpones
agent preparation when an execution-library opportunity is absent. An invalid
setting produces a notice and retains the deferred behavior. The setting is
read once per worker process, making each experimental condition explicit and
stable throughout a trial.

Optional attempt tracing records worker receipt, first preparation start,
preparation readiness, execution start, and completion. Collection is disabled
by default and applies only to auxiliary-payload tasks. The additional process
fields are generic timestamps; DataVine-specific selection and logging remain
in the worker integration point. This preserves the repository's module
boundary checks. Each trace is local to one worker clock. We compute local
intervals and never subtract timestamps from different hosts to infer latency.

For a completed attempt with timestamps $t_r,t_p,t_d,t_e,t_f$, the validation
checks $0<t_r\leq t_p\leq t_d\leq t_e\leq t_f$. The intervals $t_p-t_r$,
$t_d-t_p$, $t_e-t_d$, and $t_f-t_e$ distinguish queue waiting, preparation,
post-preparation waiting, and execution. A missing trace is not interpreted as
zero waiting. Instrumented functional runs verify one trace per physical
attempt; uninstrumented campaigns provide the primary performance samples.

\subsection{A stronger concurrency baseline}
Stock TaskVine fork libraries reject configurations with more function slots
than library cores. The one-slot-per-core baseline retains that behavior.
To test whether a speedup can be explained merely by a larger fixed window,
we also construct an experimental TaskVine control with two or four slots per
allocated core. This control changes only the startup restriction in a private,
generated library file. It does not alter repository Poncho sources, advertise
extra physical CPUs, or increase the worker's actual allocation. The worker's
slot accounting still determines admitted function calls; native math libraries
are limited to one thread.

We label these configurations as extended TaskVine throughout the evaluation.
They are not claimed to be supported stock configurations or a reimplementation
of DataVine's feedback. Their purpose is to challenge attribution: DataVine
should be compared against a baseline capable of exposing more concurrency,
not just against an artificially restrictive one. The generated source hash
and the exact transformation are retained with the trial evidence.

\subsection{Controlled dependency initialization}
A generated DataVine executor prototype preloads an explicit module list before
its readiness handshake. It is selected through the existing executor-path
configuration and leaves the default executable intact. The artifact records
the base executable hash, generated hash, and imported modules. The matching
TaskVine library hoists exactly those modules. Preloading is a deployment
control, not a new scheduling algorithm or a claim that arbitrary initialized
libraries can safely be inherited across fork. Native numerical libraries are
limited to one thread in the measured configuration.

\subsection{Recall destination progress}
The initial remote experiments exposed a liveness defect in dense least-loaded
dispatch. A successfully recalled task retained a reserved destination, while
reoffering other lower-load workers allowed the same incompatible worker to be
selected repeatedly within a bounded pass. A ready task could consequently
remain unassigned after all other work drained. The diagnostic trace showed
one such task still READY with idle execution libraries.

The fix preserves both destination binding and least-loaded preference. At the
start of each pass, the worker pool compacts stale generation entries and clears
visit marks. Selection considers each live offered worker at most once during
that pass. Reoffering an incompatible worker preserves future availability but
does not hide another worker in the current pass. This adds no per-task worker
list allocation. A deterministic regression covers repeated offers, unequal
loads, and recycled worker slots. The original four-worker diagnostic timed
out before the fix and completed with all results validated after the fix.
Pre-fix campaign data is retained separately and excluded from the primary
performance figures.

\subsection{Correctness and compatibility checks}
The full DataVine regression suite passes \UpgradeRegressionPassed{} gates. Focused
runtime acceptance executes the same 32-task phase graph under eager and
deferred preparation, verifies every logical result and physical attempt
count, and checks all trace orderings. The existing generic TaskVine serverless
test also passes in a private scratch directory; it covers both direct and
fork function libraries and validates returned values and file output. These
checks establish exercised compatibility, not exhaustive correctness for all
interleavings. The failure-injection experiment complements them by removing
only worker jobs owned by its own factory.
